\subsection{Bobinas con núcleo metálico}

Para esta parte del análisis se decidió reemplazar la bobina utilizada en los pasos anteriores, esquematizada como \(L_1\), por una segunda bobina, \(L_2\). Esta decisión fue tomada en conjunto con los profesores, debido a que los resultados medidos con \(L_1\) presentaban una diferencia significativa respecto del comportamiento esperado para el análisis propuesto en este trabajo. Por lo tanto, se consideró que existían efectos que escapaban al modelo teórico utilizado y que impedían estudiar correctamente los fenómenos principales de la experiencia.

\subsection{Diferencias entre los núcleos macizos y laminados}

Las corrientes de Foucault aparecen en los núcleos de las bobinas cuando se trabaja con circuitos de corriente alterna. Estas corrientes son corrientes inducidas que circulan en bucles cerrados dentro del material conductor del núcleo. Como consecuencia, generan campos magnéticos que se oponen a la variación del campo aplicado por la fuente de tensión. Este fenómeno provoca una disipación de energía en forma de calor, es decir, una pérdida de potencia activa.

La intensidad de estas corrientes depende tanto del tipo de material del núcleo como de su forma constructiva. Por este motivo, aun utilizando un mismo material, el diseño del núcleo puede modificar significativamente las pérdidas asociadas a las corrientes de Foucault.

Un núcleo laminado se fabrica ensamblando múltiples láminas delgadas, recubiertas o separadas entre sí mediante una capa aislante. Esta construcción interrumpe los caminos conductores de gran tamaño y obliga a las corrientes de Foucault a circular en bucles más pequeños dentro de cada lámina. Como resultado, aumenta la resistencia efectiva de dichos caminos, disminuye la magnitud de las corrientes inducidas y se reducen las pérdidas de energía y la generación de calor no deseada.

En cambio, en un núcleo macizo las corrientes de Foucault pueden circular con mayor libertad a través de una sección más grande del material. Esto permite la formación de bucles de corriente más amplios, lo que incrementa la corriente inducida y, por consiguiente, aumenta la potencia activa disipada por este efecto.

Por esta razón, los núcleos laminados son preferibles en aplicaciones de corriente alterna, ya que permiten reducir en mayor medida las pérdidas por corrientes de Foucault respecto de un núcleo macizo.
